% !TEX program = xelatex
\documentclass[11pt,a4paper]{article}

\usepackage[UTF8]{ctex}
\usepackage[margin=2.2cm]{geometry}
\usepackage{amsmath, amssymb, amsthm}
\usepackage{listings}
\usepackage{xcolor}
\usepackage{tikz}
\usepackage{booktabs}
\usepackage{multirow}
\usepackage{hyperref}
\usetikzlibrary{positioning, arrows.meta, shapes.geometric, fit, calc, decorations.pathreplacing}

\hypersetup{
  colorlinks=true,
  linkcolor=blue!60!black,
  urlcolor=teal!70!black,
}

\lstdefinelanguage{Rust}{
  morekeywords={fn,let,mut,pub,struct,impl,use,mod,self,Self,Option,Some,None,return,if,else,for,while,loop,match,break,continue,as,where,trait,type,enum,move,ref,const,static,unsafe,extern,crate,super,dyn,async,await,box,in},
  sensitive=true,
  morecomment=[l]{//},
  morecomment=[s]{/*}{*/},
  morestring=[b]",
}

\lstset{
  language=Rust,
  basicstyle=\ttfamily\footnotesize,
  keywordstyle=\color{blue!60!black}\bfseries,
  commentstyle=\color{green!40!black}\itshape,
  stringstyle=\color{red!60!black},
  numbers=left,
  numberstyle=\tiny\color{gray},
  numbersep=8pt,
  frame=single,
  framerule=0.4pt,
  rulecolor=\color{gray!40},
  breaklines=true,
  breakatwhitespace=true,
  showstringspaces=false,
  tabsize=4,
  columns=fullflexible,
  keepspaces=true,
}

\title{\textbf{halfedge\,: 顶点焊接模块设计文档} \\ \large \texttt{src/weld.rs}}
\author{halfedge 库}
\date{2026-07-04}

\begin{document}
\maketitle
\tableofcontents

\section{模块定位}

\texttt{weld.rs} 提供 \texttt{weld\_vertices}：按欧氏距离阈值合并「位置相同」
的顶点，常用于修复 OBJ/PLY 载入时因浮点精度产生的「裂缝」——同一物理顶点
被表示为多个微小偏差的副本。

\subsection{API}

\begin{center}
\renewcommand{\arraystretch}{1.18}
\begin{tabular}{@{}ll@{}}
  \toprule
  \textbf{签名} & \texttt{pub fn weld\_vertices(mesh: \&mut MeshStorage, epsilon: f64)} \\
                & \hspace*{4em}\texttt{-> Result<usize, TopologyError>} \\
  \midrule
  \textbf{输入} & \texttt{mesh}：可变借用；\texttt{epsilon}：合并阈值（$\ge 0$） \\
  \textbf{输出} & 成功时返回被移除的顶点数量 \\
  \textbf{错误} & 修复后拓扑校验失败时返回 \texttt{TopologyError} \\
  \bottomrule
\end{tabular}
\end{center}

\section{算法}

\subsection{整体流程}

\begin{center}
\resizebox{\textwidth}{!}{%
\begin{tikzpicture}[
  node distance=0.5cm,
  proc/.style={draw, rounded corners=2pt, minimum width=11cm, minimum height=0.7cm, align=center, font=\small},
  dec/.style={diamond, draw, aspect=2.4, fill=orange!12, inner sep=1pt, font=\small, align=center, minimum width=3.5cm},
  op/.style={-Stealth, thick},
  startend/.style={draw, rounded corners=10pt, minimum width=2.6cm, minimum height=0.7cm, font=\small, fill=gray!12}
]
  \node[startend] (s) {开始};
  \node[dec, below=of s] (d0) {$V < 2$？};
  \node[proc, right=1.6cm of d0, fill=green!12] (early) {返回 \texttt{Ok(0)}};
  \node[proc, below=of d0, fill=blue!8] (p1) {1. 收集所有 (VertexId, position) 到 \texttt{Vec}};
  \node[proc, below=of p1, fill=blue!8] (p2) {2. 按 $x$ 坐标排序（\texttt{sort\_by}）};
  \node[proc, below=of p2, fill=yellow!15] (p3) {3. 双重循环 $i < j$：若 $|x_j - x_i| > \epsilon$ 则 \texttt{break}，否则计算 $\|p_j - p_i\|^2 < \epsilon^2$ 时合并};
  \node[proc, below=of p3, fill=yellow!15] (p4) {4. Union-Find：\texttt{parent[rj] $\gets$ ri}（路径压缩）};
  \node[proc, below=of p4, fill=blue!8] (p5) {5. 构建 \texttt{HashMap<VertexId, VertexId>}（old $\to$ rep）};
  \node[proc, below=of p5, fill=blue!8] (p6) {6. 重映射所有半边的 \texttt{vertex} 引用};
  \node[proc, below=of p6, fill=blue!8] (p7) {7. 删除非代表顶点（\texttt{remove\_vertex}）};
  \node[proc, below=of p7, fill=red!12] (p8) {8. 清理退化面（去重后顶点数 $< 3$ 的面）};
  \node[dec, below=of p8] (d1) {剩余面数 $> 0$？};
  \node[proc, right=1.6cm of d1, fill=yellow!15] (val) {\texttt{validate\_mesh}；失败 $\to$ 返回 \texttt{Err}};
  \node[startend, below=of d1, fill=green!12] (e) {返回 \texttt{Ok(removed)}};
  \draw[op] (s) -- (d0);
  \draw[op] (d0) -- node[above, font=\scriptsize] {是} (early);
  \draw[op] (d0) -- node[right, font=\scriptsize] {否} (p1);
  \draw[op] (p1) -- (p2);
  \draw[op] (p2) -- (p3);
  \draw[op] (p3) -- (p4);
  \draw[op] (p4) -- (p5);
  \draw[op] (p5) -- (p6);
  \draw[op] (p6) -- (p7);
  \draw[op] (p7) -- (p8);
  \draw[op] (p8) -- (d1);
  \draw[op] (d1) -- node[above, font=\scriptsize] {是} (val);
  \draw[op] (d1) -- node[right, font=\scriptsize] {否} (e);
  \draw[op] (val) |- (e);
\end{tikzpicture}
}
\end{center}

\subsection{排序 + 双重循环：$O(n \log n + nk)$}

直接两两比较所有顶点对是 $O(n^2)$，对大网格不可接受。本实现按 $x$ 坐标排序后，
内层循环遇到 $|x_j - x_i| > \epsilon$ 即 \texttt{break}，将平均复杂度降至
$O(n \log n + nk)$，其中 $k$ 是 $x$ 邻域内的平均顶点数（对均匀分布网格 $k \approx$ 常数）。

\subsection{Union-Find 等价类}

合并操作使用并查集（Union-Find）：

\begin{lstlisting}
let mut parent: Vec<usize> = (0..n).collect();
// find with path halving
let mut root = i;
while parent[root] != root { root = parent[root]; }
// union: parent[rj] = ri  (ri 是 i 的根, rj 是 j 的根)
if ri != rj { parent[rj] = ri; }
\end{lstlisting}

\textbf{代表元}：等价类中 $x$ 最小的顶点（因为外层循环 $i$ 升序，
首次合并时 $i < j$，\texttt{parent[rj] = ri} 保证代表元是序号最小的）。

\subsection{重映射半边 vertex 引用}

收集 \texttt{mapping: HashMap<VertexId, VertexId>}（old $\to$ rep）后，
遍历所有半边：

\begin{lstlisting}
for he in mesh.halfedge_ids().collect::<Vec<_>>() {
    if let Some(h) = mesh.get_halfedge_mut(he) {
        if let Some(&new_v) = mapping.get(&h.vertex) {
            h.vertex = new_v;
        }
    }
}
\end{lstlisting}

\subsection{清理退化面}

合并后，可能产生顶点重复的退化面（如原 $[a, b, c]$ 中 $a, b$ 被合并为 $a$，
则面变为 $[a, a, c]$，唯一顶点数 $= 2 < 3$）。对这些面：

\begin{enumerate}
  \item 收集面的所有半边；
  \item 逐条 \texttt{remove\_halfedge}（同时解除 twin 关系）；
  \item \texttt{remove\_face}。
\end{enumerate}

\subsection{最终校验}

调用 \texttt{topology\_ops::validate\_mesh} 检查修复后网格的拓扑合法性
（twin 互指、next/prev 一致、面边界环闭合等）。失败时返回 \texttt{Err}。

\section{复杂度}

\begin{center}
\renewcommand{\arraystretch}{1.18}
\begin{tabular}{@{}lll@{}}
  \toprule
  \textbf{阶段} & \textbf{时间复杂度} & \textbf{备注} \\
  \midrule
  排序            & $O(n \log n)$              & $n = V$ \\
  双重循环        & $O(n k)$                   & $k$ = $x$ 邻域平均顶点数 \\
  Union-Find     & $O(n \alpha(n))$           & $\alpha$ = 反 Ackermann \\
  重映射半边     & $O(E)$                     & 遍历所有半边 \\
  删除冗余顶点   & $O(V)$                     & \\
  清理退化面     & $O(F \cdot \bar{k})$       & $\bar{k}$ = 平均面边数 \\
  validate\_mesh & $O(V + E + F)$             & 线性校验 \\
  \midrule
  \textbf{总计}  & $O(n \log n + n k + E)$    & 通常 $k$ 很小，接近 $O(n \log n)$ \\
  \bottomrule
\end{tabular}
\end{center}

\section{退化守卫}

\begin{itemize}
  \item $V < 2$ 时直接返回 \texttt{Ok(0)}，跳过所有处理；
  \item \texttt{partial\_cmp} 返回 \texttt{None}（NaN）时使用 \texttt{Equal}
        避免 panic；
  \item \texttt{remove\_vertex} 内部解除该顶点的所有半边引用，安全；
  \item \texttt{contains\_halfedge} / \texttt{contains\_face} 防止重复删除；
  \item 修复后 \texttt{validate\_mesh} 兜底，确保拓扑合法。
\end{itemize}

\section{测试覆盖}

\begin{center}
\renewcommand{\arraystretch}{1.18}
\begin{tabular}{@{}ll@{}}
  \toprule
  \textbf{测试名} & \textbf{验证点} \\
  \midrule
  \texttt{weld\_no\_vertices}        & 空网格返回 \texttt{Ok(0)} \\
  \texttt{weld\_single\_vertex}      & 单顶点返回 \texttt{Ok(0)} \\
  \texttt{weld\_two\_close\_vertices} & 距离 $0.001$ 的两顶点在 $\epsilon = 0.01$ 下合并为 1，第三个顶点保留 \\
  \bottomrule
\end{tabular}
\end{center}

全模块共 3 个单元测试，覆盖空网格、单顶点、合并三个核心场景；
项目整体 \texttt{cargo test} 共 371 个用例。
\textbf{待补充}：退化面清理测试、大规模网格性能测试、$\epsilon = 0$ 严格相等测试。

\section{设计权衡}

\begin{itemize}
  \item \textbf{排序 + 双重循环 vs 空间哈希}：当前实现简单稳健，对均匀分布
        顶点接近 $O(n \log n)$；若顶点高度聚集（$k$ 大），可改用空间哈希
        （uniform grid）将 $k$ 限制为常数；
  \item \textbf{Union-Find vs 直接覆盖}：并查集支持\textbf{传递合并}
        （$a \sim b, b \sim c \Rightarrow a \sim c$），避免「链式合并」错误，
        而直接覆盖可能丢失等价类信息；
  \item \textbf{仅按 $x$ 排序}：对极端情况（所有顶点 $x$ 相同但 $y, z$ 不同）
        退化为 $O(n^2)$；可改为按最长轴排序，但实现复杂度增加；
  \item \textbf{代表元选择}：选择 $x$ 最小的顶点（序号最小），保证幂等性
        （多次焊接结果相同）；不选择「最近邻」中心点，因为需要计算新位置，
        增加复杂度；
  \item \textbf{不更新 Vertex.halfedge 引用}：\texttt{remove\_vertex} 内部
        已处理解除引用，剩余顶点的入口字段保持有效；
        若代表元 vertex 的 halfedge 引用指向被删除的顶点相关半边，
        \texttt{validate\_mesh} 会报告。
\end{itemize}

\end{document}
